\section*{Chapitre \ref{ch:qm:transforms-interpolation-and-algorithmic-differentiation} --- Transformées, interpolation et différentiation algorithmique}

\begin{solution}{exo:qm:transforms-interpolation-and-algorithmic-differentiation:1}
Directement $N^2 = 16\,777\,216$~; en base~2, environ $\frac N2\log_2N = 2\,048 \times 12 = 24\,576$, soit environ 700 fois moins.
\end{solution}

\begin{solution}{exo:qm:transforms-interpolation-and-algorithmic-differentiation:2}
$(1 + \epsilon)^2 = 1 + 2\epsilon$ et $e^{1+\epsilon} = e + e\epsilon$, donc $f(1 + \epsilon) = (1 + 2\epsilon)(e + e\epsilon) = e + 3e\epsilon$~: $f'(1) = 3e \approx 8.155$.
\end{solution}

\begin{solution}{exo:qm:transforms-interpolation-and-algorithmic-differentiation:3}
$5 \cdot 2^{-k} < 10^{-10}$ impose $k > \log_2(5 \times 10^{10}) = 35.5$~: 36 itérations, comme dans le tutoriel.
\end{solution}

\begin{solution}{exo:qm:transforms-interpolation-and-algorithmic-differentiation:4}
\omterm{def:qm:transforms-interpolation-and-algorithmic-differentiation:tape}{Bande}~: $v_1 = x_1$, $v_2 = x_2$, $v_3 = v_1v_2$ (dérivées partielles $v_2$, $v_1$), $v_4 = \sin v_3$ (dérivée partielle $\cos v_3$), $v_5 = v_4 + v_1$ (dérivées partielles 1, 1). En $(1, 2)$~: $v_3 = 2$. \omterm{def:qm:transforms-interpolation-and-algorithmic-differentiation:ad}{Mode inverse}~: $\bar v_5 = 1$~; $\bar v_4 = 1$,
$\bar v_1 = 1$~; $\bar v_3 = \cos2 = -0.416$~; $\bar v_1 \mathrel{+}= \bar v_3v_2 = -0.832$, donc $\bar v_1 = 0.168$~; $\bar v_2 = \bar v_3v_1 = -0.416$. Vérification~: $\partial y/\partial x_1 = x_2\cos(x_1x_2) + 1$, $\partial y/\partial x_2 = x_1\cos(x_1x_2)$.
\end{solution}

\begin{solution}{exo:qm:transforms-interpolation-and-algorithmic-differentiation:5}
$-\ln P(t) = r(t)t$, donc $f(t) = \frac{d}{dt}(r(t)t) = r(t) + tr'(t)$. À l'extrémité longue, $t$ est grand, donc une pente de spline $r'$ de $-0.001$ par an en $t = 30$ abaisse le taux à terme de 3 points de pourcentage~: de petites ondulations de la
courbe des taux zéro deviennent de grandes ondulations du taux à terme, surtout aux longues échéances.
\end{solution}

\begin{solution}{exo:qm:transforms-interpolation-and-algorithmic-differentiation:6}
Le payoff $\mathbf 1\{S_T > K\}$ est constant par morceaux le long de la trajectoire, donc sa dérivée trajectorielle est nulle presque sûrement~: dériver chaque payoff simulé rate le saut. Remèdes~: lisser le payoff (un
call spread étroit), utiliser la méthode du rapport de vraisemblance pour ce terme, intégrer analytiquement le dernier pas (\omterm{def:qm:probability-at-speed:condexp}{espérance conditionnelle}), ou combiner les deux (Vibrato Monte Carlo).
\end{solution}

\begin{solution}{exo:qm:transforms-interpolation-and-algorithmic-differentiation:7}
\texttt{implied\_vol\_iterations()}~: dichotomie 36 à chaque strike~; sécante 10, 7, 4, 7, 18 et échec~; Newton 9, 6, 3, 6, 13 et échec~; Brent 17, 17, 5, 17, 17 et 18. À 400, le prix
($2.5 \times 10^{-7}$) et le vega au point de départ ($8 \times 10^{-9}$) sont minuscules~: le premier pas de Newton part trop loin, jusqu'à une volatilité de 30.6, où le vega subit un sous-dépassement, et la pente de la sécante est tout aussi
peu fiable~; les méthodes d'encadrement ne peuvent pas sortir de $[0.001, 5]$.
\end{solution}

\begin{solution}{exo:qm:transforms-interpolation-and-algorithmic-differentiation:8}
Une différence centrée a une \omterm{def:qm:finite-difference-methods:consistent}{erreur de troncature} d'ordre $h^2f^{\prime\prime\prime}$ et une erreur d'arrondi d'ordre $u|f|/h$~; à l'échelle du livre (valeur d'environ 112, dérivées seconde et troisième grandes au voisinage des strikes des caplets), un
écart de $10^{-5}$ sur 500 représente une erreur relative de $2 \times 10^{-8}$, soit la précision de la perturbation elle-même. L'adjoint concorde avec le \omterm{def:qm:transforms-interpolation-and-algorithmic-differentiation:ad}{mode direct} à $10^{-13}$ près~: la perturbation est la moins précise des deux.
\end{solution}

\begin{solution}{pb:qm:transforms-interpolation-and-algorithmic-differentiation:1}
\textbf{1.} Décroissance exponentielle~: Black--Scholes exact à $7 \times 10^{-15}$ près avec 64 termes~; Merton $10^{-7}$ à 64 termes et $2 \times 10^{-12}$ à 128.
\textbf{2.} $\P(S_T > 110) = 0.3503999$, à comparer au mélange poissonnien de lognormales à $3 \times 10^{-10}$ près.
\textbf{3.} Face à 400\,000 trajectoires simulées~: 9.5243 contre $9.528 \pm 0.022$.
\textbf{4.} Le payoff du put est borné, donc ses coefficients en cosinus restent petits et la troncature de la queue supérieure n'amplifie pas les erreurs~; la parité call-put donne le call.
\textbf{5.} Lorsqu'un ensemble complet de strikes est nécessaire d'un coup sur une grille régulière, comme dans la méthode de Carr et Madan.
\textbf{6.} Ils tombent à 2.63~\% à trente ans~: la condition de fin naturelle laisse la pente du spline chuter fortement, et $tr'(t)$ amplifie la pente aux longues échéances.
\textbf{7.} $rt$ monotone, donc des taux à terme non négatifs lorsque les données le permettent~; ici chaque taux à terme se situe entre 3.32~\% et 5.67~\%.
\textbf{8.} Ses taux à terme sautent à chaque pilier, de 0.43 point de pourcentage au plus.
\textbf{9.} À la monnaie, Newton 3, sécante 4, Brent 5, dichotomie 36~; à 400, Newton et la sécante échouent, Brent nécessite 18 itérations et la dichotomie 36.
\textbf{10.} 1.92 près des extrémités.
\textbf{11.} 2\,864 opérations élémentaires et 3\,698 couples (opération, entrée).
\textbf{12.} La perturbation exige 401 évaluations, le \omterm{def:qm:transforms-interpolation-and-algorithmic-differentiation:ad}{mode direct} 1\,433, le \omterm{def:qm:transforms-interpolation-and-algorithmic-differentiation:ad}{mode inverse} 3.58 selon le nombre d'opérations.
\textbf{13.} Les \omterm{def:qm:transforms-interpolation-and-algorithmic-differentiation:ad}{modes inverse} et direct concordent à $2 \times 10^{-13}$ près~; le \omterm{def:qm:transforms-interpolation-and-algorithmic-differentiation:ad}{mode inverse} et les perturbations centrées concordent à $9 \times 10^{-6}$ près, soit l'\omterm{def:qm:finite-difference-methods:consistent}{erreur de troncature} et d'arrondi des perturbations.
\textbf{14.} Avec une interpolation linéaire, chaque date de flux dépend des deux piliers qui l'encadrent~; les piliers sans date dans leur voisinage ont une sensibilité nulle.
\textbf{15.} Il conserve au plus 201 états au lieu de 10\,000 (60\,003 nœuds de \omterm{def:qm:transforms-interpolation-and-algorithmic-differentiation:tape}{bande}), au prix d'une passe avant supplémentaire, avec le même gradient à $10^{-14}$ près.
\textbf{16.} Sensibilités des payoffs discontinus et des entrées qui passent par du code non dérivable (recherches dans des tables, branchements), et un échantillon nocturne en guise de contrôle.
\textbf{17.} Adjoint contre perturbation sur un échantillon d'entrées, taille et mémoire de la \omterm{def:qm:transforms-interpolation-and-algorithmic-differentiation:tape}{bande}, et vérification que la \omterm{def:qm:transforms-interpolation-and-algorithmic-differentiation:tape}{bande} était vidée à chaque valorisation.
\textbf{18.} Le coût réel, y compris le trafic mémoire de la \omterm{def:qm:transforms-interpolation-and-algorithmic-differentiation:tape}{bande}~: le décompte des opérations l'ignore~; le test en C++ mesure moins de vingt évaluations.
\textbf{19.} Résultat retenu~: \emph{un seul balayage au lieu de 401}~: le gradient adjoint du portefeuille à 400 entrées coûte 3.58 évaluations selon le nombre d'opérations, contre 401 pour la perturbation, et concorde avec le gradient perturbé à
$9 \times 10^{-6}$ près (avec le \omterm{def:qm:transforms-interpolation-and-algorithmic-differentiation:ad}{mode direct} à $2 \times 10^{-13}$ près).
\textbf{20.} La fonction risque demande les sensibilités d'un seul nombre à de nombreuses entrées, et le coût du \omterm{def:qm:transforms-interpolation-and-algorithmic-differentiation:ad}{mode inverse} ne dépend pas du nombre de ces entrées.
\end{solution}

\begin{solution}{iq:qm:transforms-interpolation-and-algorithmic-differentiation:1}
Le \omterm{def:qm:transforms-interpolation-and-algorithmic-differentiation:ad}{mode direct} propage les dérivées le long d'une direction d'entrée, en même temps que les valeurs~: coût proportionnel au nombre d'entrées, avantageux pour peu d'entrées et beaucoup de sorties. Le \omterm{def:qm:transforms-interpolation-and-algorithmic-differentiation:ad}{mode inverse} enregistre le calcul et
propage les adjoints d'une sortie vers toutes les entrées~: coût de quelques multiples d'une évaluation, indépendant du nombre d'entrées, avantageux pour une sortie et beaucoup d'entrées, au prix de la mémoire.
\iqlookfor{Croissance du coût avec le nombre d'entrées ou de sorties, et mémoire de la \omterm{def:qm:transforms-interpolation-and-algorithmic-differentiation:tape}{bande}.}
\end{solution}

\begin{solution}{iq:qm:transforms-interpolation-and-algorithmic-differentiation:2}
Un type actif dont les opérateurs enregistrent (entrées, dérivées partielles) sur une \omterm{def:qm:transforms-interpolation-and-algorithmic-differentiation:tape}{bande}, et une boucle inverse parcourant la \omterm{def:qm:transforms-interpolation-and-algorithmic-differentiation:tape}{bande} qui accumule les adjoints. Limites~: mémoire de la \omterm{def:qm:transforms-interpolation-and-algorithmic-differentiation:tape}{bande} (\omterm{def:qm:transforms-interpolation-and-algorithmic-differentiation:tape}{checkpointing}), sûreté des threads (une \omterm{def:qm:transforms-interpolation-and-algorithmic-differentiation:tape}{bande} par
fil d'exécution), code tiers non templatisé, branchements et points anguleux, et surcoût de l'enregistrement des boucles internes, que réduisent les templates d'expressions ou les adjoints écrits à la main pour les noyaux.
\iqlookfor{Conception de la \omterm{def:qm:transforms-interpolation-and-algorithmic-differentiation:tape}{bande}, mémoire et au moins deux limites en production.}
\end{solution}

\begin{solution}{iq:qm:transforms-interpolation-and-algorithmic-differentiation:3}
Encadrer la volatilité (le prix est croissant en elle entre des bornes connues), utiliser une méthode sécurisée telle que la \omterm{def:qm:transforms-interpolation-and-algorithmic-differentiation:roots}{méthode de Brent}, ou Newton avec repli sur la dichotomie, en partant d'une bonne
approximation rationnelle~; travailler avec des prix normalisés, et traiter comme des erreurs les prix inférieurs à la valeur intrinsèque ou supérieurs à la borne supérieure.
\iqlookfor{Encadrement avec un pas intérieur rapide, et contrôles de bornes.}
\end{solution}

\begin{solution}{iq:qm:transforms-interpolation-and-algorithmic-differentiation:4}
Interpoler une grandeur dont la forme porte l'économie, comme $\ln P(t)$ ou $rt$, avec un schéma local monotone, de sorte que les taux à terme restent positifs, continus là où on le souhaite, et locaux~: une variation de cotation
ne déplace la courbe qu'au voisinage de son pilier. Une \omterm{def:qm:transforms-interpolation-and-algorithmic-differentiation:spline}{spline cubique naturelle} sur les taux zéro est globale (une seule cotation déplace tous les taux à terme) et peut produire des taux à terme oscillants ou négatifs.
\iqlookfor{Localité et courbe à terme.}
\end{solution}

\begin{solution}{iq:qm:transforms-interpolation-and-algorithmic-differentiation:5}
Par des méthodes de Fourier~: le développement COS ou la transformée amortie de Carr et Madan avec une \omterm{def:qm:transforms-interpolation-and-algorithmic-differentiation:dft}{transformée de Fourier rapide}, ou une intégrale de Gil-Pelaez pour les probabilités~; vérifier l'intervalle de troncature à partir des \omterm{def:qm:jump-processes:cumulant}{cumulants}, et
valider par simulation.
\iqlookfor{COS ou FFT, avec un choix de troncature et un contrôle.}
\end{solution}

\begin{solution}{iq:qm:transforms-interpolation-and-algorithmic-differentiation:6}
L'indicateur de barrière est constant par morceaux le long de chaque trajectoire, donc la dérivée trajectorielle manque la discontinuité. Utiliser une barrière lissée (une probabilité de survie par pas du \omterm{def:qm:brownian-motion:bridge}{pont brownien}, qui est
dérivable), un terme de rapport de vraisemblance, ou des \omterm{def:qm:probability-at-speed:condexp}{espérances conditionnelles} à la barrière~; puis vérifier contre un prix perturbé, avec des nombres aléatoires communs.
\iqlookfor{Discontinuité et reformulation dérivable.}
\end{solution}
